\section{Error Analysis}
\label{sec:error-analysis}

\subsection{Checkpoint-bound threshold transfer}
\label{sec:threshold-transfer}

The primary operating point is selected on the training-time dev subset and
is stored with the exact best-dev checkpoint that produced it.  Held-out
predictions use that pair unchanged.  To describe calibration transfer
without tuning on test, we may compute an oracle diagnostic after primary
scoring,
\begin{equation}
  \Delta_{\mathrm{cal}} =
  F1_{\mathrm{test}}(\tau^{\star}_{\mathrm{test}})
  -F1_{\mathrm{test}}(\tau_{\mathrm{dev}}),
  \label{eq:calibration-gap}
\end{equation}
where $\tau^{\star}_{\mathrm{test}}$ is used only to quantify the possible
threshold-transfer gap.  It never replaces $\tau_{\mathrm{dev}}$, changes a
reported primary metric, or influences the once-only final evaluation.  We
also inspect precision and recall at $\tau_{\mathrm{dev}}$ to distinguish a
ranking failure from a shifted score scale.

\ifdiagnosticdraft
Calibration-gap values are intentionally blank until corrected H100
predictions are available.  Historical V100 thresholds are not eligible for
this analysis because they are tied to the superseded selection contract.
\fi

\subsection{Dark vessels and near-shore clutter}
\label{sec:slice-errors}

Dark-vessel recall and near-shore F1 probe different failure modes.  The
former measures misses among human-verified vessels without a matched AIS
report; the latter includes unmatched predictions near the coast and
therefore exposes shoreline clutter.  Dark labels occur only in the
50-scene human-verified set.  Conversely, the 16-scene test split contains
only two valid near-shore vessel targets, too little support for a stable
standalone coastal conclusion.  We therefore show test near-shore F1 only as
a support-qualified diagnostic and reserve substantive slice interpretation
for the once-only human-verified evaluation in
Fig.~\ref{fig:final-slices}.

When the final predictions are unlocked, we compare slice errors at each
model's frozen dev threshold.  The analysis reports the numerator and
denominator beside every recall or F1 value so that an apparently large gap
cannot conceal small support.  A source-domain advantage in aggregate F1 is
not assumed to transfer to dark vessels or coastal clutter; those are
separate empirical questions.

\subsection{False-positive taxonomy and qualitative review}
\label{sec:qualitative-errors}

Figure~\ref{fig:qualitative} presents registered $[\mathrm{VH},\mathrm{VV},
\mathrm{VH}-\mathrm{VV}]$ chips with ground truth, predictions, confidence,
distance to shore, and the frozen operating threshold.  A stratified manual
review of approximately 200 false positives assigns one of four
predeclared descriptive categories: shoreline clutter, fixed infrastructure,
open-water sea clutter, or sidelobe/bright-target response.  Ambiguous cases
remain marked ambiguous rather than being forced into a class.  The sample is
used to explain mechanisms, not to estimate population frequencies or to
retune the detector.

\ifdiagnosticdraft
The diagnostic build contains the gallery layout and category definitions
only.  Image examples and category counts must come from the corrected H100
cohort and the authorized evaluation stage before they can enter the
submission.
\else
The qualitative cases are interpreted alongside, but do not supersede, the
quantitative findings in Sections~\ref{sec:transfer-results}
and~\ref{sec:final-results}.
\fi
